\section{Introduction}

The measurement problem is usually introduced as a tension among unitary dynamics, the
superposition principle, and the apparent definiteness of observed outcomes. On that framing, the
central question is how a linear, non-collapsing state description yields a world in which observers
report definite results. Collapse theories answer by changing dynamics, hidden-variable theories by
adding ontological structure, Everettian theories by retaining unitary dynamics while reinterpreting
branching, and instrumentalist or QBist views by weakening the descriptive burden placed on the
theory. The debate is then conducted primarily as a dispute over outcome production.\footnote{For
standard landmarks see \citet{epr1935}, \citet{bell1964}, \citet{zurek2003},
\citet{schlosshauer2005}, and \citet{wallace2012}.}

This manuscript isolates a prior problem. Outcome reports enter scientific practice through
measurement records. Those records are not inert marks. They anchor claims about what was
observed, what is now remembered, what may be anticipated, and how rival hypotheses are to be
compared. If a theory is used to underwrite such discourse, then the theory must stabilize reference to
those records strongly enough for their evidential role to be fixed. Before asking how one outcome is
produced, one must therefore ask what makes a measurement record the determinate referent of the
discourse in which it is used.

That question becomes acute in physically non-selective theories. In Everettian quantum mechanics,
for example, unitary dynamics and decoherence can stabilize multiple post-measurement
branch-histories, each containing a determinate branch-local observer and a determinate branch-local
record, without providing a physical selector that privileges one such history as \emph{the} referent of
pre-measurement diachronic discourse. Branch-local determinacy is therefore not yet diachronic
referential determinacy. The issue is not whether post-branch observers have determinate experiences.
The issue is whether the theory fixes which future record is the referent of pre-measurement claims of
the form \enquote{I will observe $o$}, \enquote{I remember observing $o$}, or \enquote{this record
confirms $H_1$ over $H_2$}.

The argument of the paper is conditional but exact. Work throughout on a fixed downstream regime in
which admissibility is already forced, standing is bivalent on the evaluated fragment, same-scope
admissibility classification is unique, the admissible interior is unique up to lawful equivalence, and
standing is conserved. Within that fixed regime, define admissible reference assignments for
diachronic measurement discourse and consider the standing-bearing content they induce. The central
mathematical claim is then a \emph{measurement-reference quotient theorem}: any admissibility-
preserving presentation of the reference-bearing content of measurement-record discourse factors
uniquely through a quotient of admissible assignments by standing-equivalence. Differences that do
not survive that quotient are semantic skin. Differences that do survive it are standing-relevant. The
next hardening step is an \emph{equivalence-collapse theorem}: any same-scope quotient relation that
claims to preserve standing-bearing content collapses to that same quotient. In particular, divergence
that changes confirmational role is not a harmless repackaging of one content; it is quotient-relevant
divergence.

Once that quotient structure is in place, the rest of the theorem ladder tightens around one bridge.
First, quotient uniqueness and global equivalence collapse show that same-scope standing-bearing
content is fixed by the measurement-reference quotient rather than by a discretionary presentation.
Second, fixation commutation is identified as the bridge invariant: if lawful redescription does not
preserve the relevant quotient class, then the theory fails to preserve its own evidential target. Third,
the bare Everettian package is shown to be physically non-selective and therefore not to supply that
bridge by dynamics alone. Fourth, theory-level confirmation is proved to be undefined, not merely
weakened, whenever fixation fails. Fifth, memory, anticipation, deliberation, and agency inherit the
same dependence because any admissible diachronic claim factors through the same quotient-level
lineage. Finally, interpretation packages are exhausted by the structural locus at which singular
reference is fixed, if it is fixed at all; no fifth same-scope locus remains.

Two restrictions on scope matter throughout. First, the paper is not a general semantics of quantum
language. It concerns a narrower target practice: truth-apt diachronic discourse in which
measurement records function as evidential anchors for memory, anticipation, deliberation, and
theory-level confirmation. Second, the paper does not re-derive the downstream closure package on
which it works. That package is imported explicitly and cited where needed; the present manuscript
uses it to sharpen the measurement argument rather than to reproduce the companion closure corpus.

The resulting diagnosis is different from the standard presentation of the measurement problem. The
first failure is not yet that the theory does not produce a unique outcome. The first failure is that, in a
physically non-selective setting, the theory may leave underfixed the referents of the very records
through which outcomes are reported and used as evidence. On that target discourse, the measurement
problem is therefore a problem of reference to measurement records before it is a problem of outcome
production.

\begin{claimbox}
\textbf{Central thesis.} Let $T$ be a physically non-selective theory used to underwrite truth-apt
diachronic discourse about measurement, memory, anticipation, deliberation, and theory-level
confirmation. Then admissible reference assignments for measurement records must either collapse into
a single standing-equivalence class on that discourse or else make confirmation and diachronic
practice assignment-relative rather than theory-relative. Bare Everettian dynamics, entanglement,
decoherence, and branch-local record stability do not by themselves collapse the admissible class.
Accordingly, the measurement problem is a problem of referential determinacy for measurement
records, not merely a problem of physical outcome production.
\end{claimbox}

The proof order is strict. Section~\ref{sec:targetclass} states the fixed downstream substrate, gives a
compressed substrate-necessity bridge, and fixes the exact standing conditional. Section~\ref{sec:semanticcore}
defines admissible reference assignments in gate-governed fixed-domain form, proves the
measurement-reference quotient theorem, proves the global equivalence-collapse theorem that blocks
the charge that the quotient was chosen rather than forced, and then promotes fixation commutation
to the bridge invariant between quotient structure and evidential use. Section~\ref{sec:branching}
proves the Everettian non-selection results as the physical trigger for the bridge problem.
Section~\ref{sec:confirmation} proves the theory-level confirmation criterion and the resulting
stabilization consequences. Section~\ref{sec:toy} exhibits a worked reversal witness.
Section~\ref{sec:diachronic_practice} proves the diachronic closure and inheritance theorems for
memory, anticipation, deliberation, and confirmation. Section~\ref{sec:taxonomy} derives selector
exhaustion from structural constraints before stating the comparative interpretation taxonomy.
Section~\ref{sec:objections} opens with a strict denial table so that standard objections are tied to
exact labeled failures.
Section~\ref{sec:nonclaims} marks the manuscript's precise limits. Appendix~\ref{app:formal-quotient}
supplies a compact formal appendix for the quotient presentation, Appendix~\ref{app:later-architecture}
states the exact companion closure package on which the main body is conditional, and
Appendix~\ref{app:audit-materials} records the dependency and objection maps.
